\needspace{4\baselineskip}
\section{Quality and Erosion Analysis}\label{sec:appendix-quality}

\subsection{Code Quality Examples}\label{sec:appendix-quality-examples}

This subsection reproduces the two \codesearch{} code fragments referenced in \autoref{sec:quality-dims}. \autoref{lst:dispatch-erosion} illustrates structural erosion, and \autoref{fig:verbosity-example} illustrates verbosity.

\begin{figure}[h]
\begin{lstlisting}[caption={Structurally eroded \texttt{find\_matches\_in\_file()} in \codesearch{} at $\ckpt{5}$ (117 lines total). Nearly all decision-point mass ends up in one function.},label={lst:dispatch-erosion}]
def find_matches_in_file(text, path, language, rules):
    ...
    for rule in rules:
        if kind == "exact":
            ...
        elif kind == "regex":
            ...
        elif kind == "pattern":
        else: ...

        if iterable is not None:
            for match in iterable:
                ...
        if kind == "pattern":
            for match in iter_pattern_matches(...):
                ...
        for node in iter_tree_nodes(source_root):
            if not node_matches_selector(selector, node):
                ...
    return matches
\end{lstlisting}
\end{figure}

\begin{figure}[h]
\begin{lstlisting}[caption={Overly verbose code from \codesearch{}. Identity list comprehension instead of \texttt{filter}, empty checks instead of building around iteration, and single-use variables.},label={fig:verbosity-example}]
for posix_path, full_path, language in source_files:
    applicable_rules = [
        r for r in all_compiled_rules
        if language in r["languages"]
    ]
    if not applicable_rules:
        continue
    with open(full_path, "r", encoding=encoding) as f:
        content = f.read()
    matches = find_matches_in_content(
        content, applicable_rules, language
    )
    if not matches:
        continue
    match_list.extend(matches)
all_matches = deduplicate_matches(match_list)
return all_matches
\end{lstlisting}
\end{figure}

\subsection{Erosion Example: \texttt{forge/main} Across Eight Checkpoints}\label{sec:appendix-erosion-example}

On \texttt{forge}, Opus~4.7's \texttt{main()} grows \erosionForgeMainCcMultiplier$\times$ in cyclomatic complexity over \erosionForgeNumCkpts{} checkpoints, from CC \erosionForgeMainCcFirst{} to CC \erosionForgeMainCcLast{}, and from \erosionForgeMainLinesFirst{} to \erosionForgeMainLinesLast{} lines. The increase is monotonic across every checkpoint. \autoref{lst:erosion-main-cp1} shows the $\ckpt{1}$ implementation, a six-line \texttt{argparse} setup followed by a small dispatch ladder. \autoref{lst:erosion-main-cp8} shows the $\ckpt{8}$ implementation, in which \texttt{argparse} has been removed and replaced with a manual \texttt{while}-loop that handles each flag in two parallel branches, one for the \verb|--X val| form and one for \verb|--X=val|. This pair of branches is repeated for seven flags before control passes to a nested resource-by-command dispatch ladder. No shared helper is introduced across the seven flag pairs, despite the per-flag template being uniform.

\begin{figure}[h]
\begin{lstlisting}[caption={\texttt{main()} from Opus~4.7's \texttt{forge} solution at $\ckpt{1}$. Argument parsing is delegated to \texttt{argparse}; dispatch is a single \texttt{if/elif} ladder.},label={lst:erosion-main-cp1}]
def main():
    parser = argparse.ArgumentParser(add_help=False)
    parser.add_argument("--data-dir", required=True)
    parser.add_argument("resource")
    parser.add_argument("command")
    parser.add_argument("rest", nargs=argparse.REMAINDER)
    args = parser.parse_args()

    if args.resource != "blueprint":
        emit_error("validation_error", f"unknown resource: {args.resource}")

    cmd, rest = args.command, args.rest
    if cmd == "create":   cmd_create(args.data_dir)
    elif cmd == "list":   cmd_list(args.data_dir)
    elif cmd == "get":    cmd_get(args.data_dir, rest[0])
    # ... two more branches ...
    else:
        emit_error("validation_error", f"unknown command: {cmd}")
\end{lstlisting}
\end{figure}

\begin{figure}[h]
\begin{lstlisting}[caption={\texttt{main()} from Opus~4.7's \texttt{forge} solution at $\ckpt{8}$, abridged to expose structure. The \texttt{argparse} call has been removed; each of seven flags is handled by an identical pair of branches. After parsing, a nested \texttt{if resource == ...: if cmd == ...} ladder dispatches over seven resources.},label={lst:erosion-main-cp8}]
def main():
    argv = sys.argv[1:]
    data_dir = revision_arg = config_path = role = scope = None
    rules_file_flag = None
    cleaned, i = [], 0
    while i < len(argv):
        tok = argv[i]
        if tok == "--data-dir":
            if i + 1 >= len(argv):
                emit_error("validation_error", "--data-dir requires a value")
            data_dir = argv[i + 1]; i += 2
        elif tok.startswith("--data-dir="):
            data_dir = tok[len("--data-dir="):]; i += 1
        # ... same two-branch pattern repeats verbatim for
        #     --revision, --config, --rules-file, --role, --scope
        #     (six more flags, twelve more branches) ...
        else:
            cleaned.append(tok); i += 1

    # ... config load, revision negotiation (~30 lines) ...

    if resource == "blueprint":
        if cmd == "create":  cmd_create(data_dir, ...)
        elif cmd == "list":  cmd_list(data_dir, ...)
        # ... five more cmd branches ...
        else: emit_error("validation_error", f"unknown command: {cmd}")
    elif resource == "revision":
        # ... same cmd-ladder shape, six more resources ...
\end{lstlisting}
\end{figure}

\subsection{Erosion Sensitivity}\label{sec:erosion-sensitivity}

We vary the high-CC cutoff (8, 10, 12) and the size term (no size term, $\sqrt{\text{SLOC}}$, linear SLOC) around the reported erosion family. Across all nine variants, the result is stable: predictive correlation with next-checkpoint pass rate stays near zero, while predictive correlation with next-checkpoint cost stays positive.

% Auto-generated by scripts/plotting/appendix_erosion_sensitivity.py
% Re-run: uv run python scripts/plotting/appendix_erosion_sensitivity.py
\begin{table}[h]
\centering
\small
\caption{Erosion metric sensitivity sweep.}
\label{tab:appendix-erosion-sensitivity}
\begin{tabular}{@{}lrr@{}}
\toprule
\textbf{Metric} & \textbf{Pass} & \textbf{Cost} \\
\midrule
Reported erosion & $-0.018$ & $0.167$ \\
Best variant (CC > 12 + Linear SLOC) & $-0.136$ & $0.095$ \\
LOC & $-0.212$ & $0.502$ \\
Mean CC & $-0.046$ & $0.218$ \\
Max CC & $-0.117$ & $0.356$ \\
Clone Ratio & $-0.042$ & $0.114$ \\
AST-grep hits / LOC & $-0.052$ & $0.151$ \\
\bottomrule
\end{tabular}
\end{table}


Size-heavy baselines such as LOC and max CC remain stronger raw cost predictors, but the erosion-family conclusion does not depend on the exact threshold or size term.
